% Proposition de fragment LaTeX pour le Chapitre 3.
% À adapter après vérification de l'exécution CI sur GitHub.
\subsection{Préparation technique de l'expérimentation}

Afin de disposer d'un support reproductible pour l'étude des pratiques XP,
nous avons retenu un kata de gestion de réunions. La première étape consiste
à préparer une mini-API FastAPI et un test automatisé de disponibilité.
Les règles métier de réservation, en particulier la validation de la
capacité d'une salle, constituent le périmètre des itérations expérimentales
à venir.

La Definition of Done technique prévoit des critères d'acceptation explicites,
des tests automatisés réussis, une revue de code et une intégration continue
avec un résultat positif. Un workflow GitHub Actions est préparé pour
exécuter les tests à chaque modification publiée ; sa réussite doit être
attestée par les journaux d'exécution du dépôt, et non supposée.

Le tableau initial du groupe comporte les colonnes \textit{Backlog},
\textit{To Do}, \textit{Doing}, \textit{Test/Review} et \textit{Done}.
Pour les itérations suivantes, nous proposons une limite WIP de deux tâches
en \textit{Doing} et de deux tâches en \textit{Test/Review}. Il s'agit
d'une règle d'amélioration proposée et non d'une limite déjà observée.
